\chapter{Sociétés de gestion, fonds de pension, fonds souverains et assureurs}\label{ch:in:asset-managers-pensions-sovereign-funds-and-insurers}

Le \omterm{def:in:asset-managers-pensions-sovereign-funds-and-insurers:swf}{fonds souverain} norvégien valait 21\,268 milliards de couronnes à la fin de 2025, soit environ 2.05~billions de dollars au
taux de change moyen de l'année, et sa société de gestion employait 678 personnes. Il a dépensé 3.8 points de base du fonds pour sa gestion, dont 39~\%
de ce montant est allé aux gestionnaires externes qui géraient 5~\% de ses actifs. Le fonds de pension national canadien, moins d'un tiers de cette taille,
a versé à ses gestionnaires externes 4.7~milliards de dollars canadiens en frais de base et commissions de performance au cours de son dernier exercice, soit 2.7 fois ses propres
charges d'exploitation. Les deux institutions ont fait des choix opposés sur ce qu'elles font elles-mêmes, et ces choix déterminent combien
de quants elles emploient et ce que font les quants. Ce chapitre décrit l'autre face des marchés~: les \omterm{def:in:asset-managers-pensions-sovereign-funds-and-insurers:owner}{propriétaires d'actifs}
dont l'argent est géré par le reste de l'industrie, les sociétés de gestion qui les servent, et les assureurs qui gèrent en
fonction de leurs passifs.

\section{Propriétaires d'actifs et leurs équipes internes}

\begin{definition}[Propriétaire d'actifs]\label{def:in:asset-managers-pensions-sovereign-funds-and-insurers:owner}
Un \emph{propriétaire d'actifs}\index{propriétaire d'actifs} est une institution qui investit les actifs qu'elle détient pour le compte de bénéficiaires ou de
l'État -- un fonds de pension, un \omterm{def:in:asset-managers-pensions-sovereign-funds-and-insurers:swf}{fonds souverain}, une dotation, le compte général d'un assureur -- et répond des
résultats devant eux, qu'il gère ces actifs lui-même ou par l'intermédiaire de gestionnaires externes.
\end{definition}

\begin{definition}[Fonds souverain, régime de retraite à prestations définies]\label{def:in:asset-managers-pensions-sovereign-funds-and-insurers:swf}
Un \emph{fonds souverain}\index{fonds souverain} est un fonds d'investissement à vocation spécifique, détenu par une
administration publique, créé à des fins macroéconomiques, qui investit (entre autres) dans des actifs financiers étrangers. Un
\emph{régime de retraite à prestations définies}\index{régime de retraite à prestations définies} promet à ses membres une prestation fixe et préétablie
versée à la retraite, de sorte que le risque d'investissement pèse sur le promoteur et que les actifs du régime soient gérés au regard des
prestations promises.
\end{definition}

Le livre~1 (chapitre~3) a présenté les \omterm{def:in:asset-managers-pensions-sovereign-funds-and-insurers:owner}{propriétaires d'actifs} comme les clients des sociétés de gestion. En tant qu'employeurs, ils se distinguent de tous les autres
acteurs de ce livre sur un point~: leur objectif est fixé par un passif ou un mandat émanant d'un gouvernement ou d'un conseil d'administration, et non par un profit pour
leurs propriétaires. Les quants d'un fonds de pension travaillent à servir les prestations~; ceux d'un \omterm{def:in:asset-managers-pensions-sovereign-funds-and-insurers:swf}{fonds souverain} visent un objectif de rendement assorti d'un budget de risque fixé
par son ministère~; ni les uns ni les autres ne touchent une part de ce qu'ils génèrent.

\begin{definition}[Gestion interne]\label{def:in:asset-managers-pensions-sovereign-funds-and-insurers:internal}
La \emph{gestion interne}\index{gestion interne} est la gestion des actifs d'un \omterm{def:in:asset-managers-pensions-sovereign-funds-and-insurers:owner}{propriétaire d'actifs} par son propre personnel
plutôt que par des gestionnaires externes au titre de mandats~; sa part est la fraction des actifs gérés de cette façon.
\end{definition}

\section{Où se situent les quants~: allocation d'actifs, risque, passifs, facteurs et indices}

Chez un \omterm{def:in:asset-managers-pensions-sovereign-funds-and-insurers:owner}{propriétaire d'actifs}, le travail quantitatif se répartit en cinq catégories.
\begin{enumerate}
\item \textbf{Allocation d'actifs}~: la répartition stratégique et son risque, la décision unique la plus importante du propriétaire.
\item \textbf{Risque}~: mesure et limites appliquées aux mandats internes et externes (livre~6).
\item \textbf{Appariement au passif}~: pour les fonds de pension et les assureurs, la gestion adossée au passif (LDI) du livre~2 (chapitre~7)~:
courbes, duration, inflation, programmes de couverture.
\item \textbf{Facteurs et actions quantitatives}~: portefeuilles systématiques internes et suivi des portefeuilles externes
(livre~8, chapitre~29).
\item \textbf{Indexation et mise en œuvre}~: réplication à faible coût des indices de référence, trading, gestion de transition.
\end{enumerate}
La combinaison dépend du choix du propriétaire entre faire et acheter. Un propriétaire qui gère la plupart de ses actifs en interne emploie
des gérants de portefeuille, des traders et des chercheurs~; un propriétaire qui achète la gestion emploie des sélectionneurs de gérants, des analystes des risques et
des responsables de l'allocation.

\section{Deux propriétaires, deux modèles}

\begin{dated}{2026-05}{Deux propriétaires d'actifs en chiffres}\label{dat:in:asset-managers-pensions-sovereign-funds-and-insurers:owners}
\textbf{Norges Bank Investment Management} (rapport annuel 2025)~: valeur du fonds~: 21\,268 milliards de NOK~; 678 salariés à Oslo,
Londres, New York et Singapour~; frais de gestion de 7\,537 millions de NOK, soit 3.8 points de base des encours sous gestion, dont
frais de personnel de 2\,476 millions de NOK, frais de base versés aux gestionnaires externes de 1\,756 millions de NOK et commissions de performance de 1\,182
millions de NOK~; 1\,062 milliards de NOK (5.0~\% du fonds) répartis sur 111 mandats externes. \textbf{CPP Investments} (exercice clos le 31
mars 2026)~: actifs nets de 793.3 milliards de dollars canadiens~; charges d'exploitation de 1\,757 millions de dollars canadiens (soit un ratio de charges d'exploitation de 23.1 points de base
points de base)~; 1\,976 millions de dollars canadiens de frais de gestion d'investissement et 2\,758 millions de dollars canadiens de commissions de performance versés aux
gestionnaires externes~; 753 millions de dollars canadiens de coûts de transaction~; investissements nets par salarié de 364 millions de dollars canadiens.
\end{dated}

Les deux chiffres que publient les propriétaires, 3.8 et 23.1 points de base, ne sont pas comparables~: l'un inclut les frais des gestionnaires externes
et l'autre non. Ramenés sur une même base à partir des totaux publiés, avec les actifs de fin d'année comme dénominateur, le
coût total de gestion du fonds norvégien s'élève à 3.5 points de base et celui du fonds canadien à 81.8, dont 59.7 sont dus aux
gestionnaires externes (voir la \cref{fig:in:asset-managers-pensions-sovereign-funds-and-insurers:costs}). La ventilation du fonds norvégien
selon le gestionnaire des actifs explique l'écart~: la part gérée en interne, qui représente 95~\%, coûte 2.3 points de base, et ses mandats externes 27.7
points de base, dont 16.5 de frais de base et le reste en commissions de performance.

\begin{omfigure}
\begin{tikzpicture}
\begin{axis}[width=10.5cm,height=5.4cm,xbar,bar width=9pt,xmin=0,xmax=90,xlabel={\small coût de gestion (points de base des actifs)},
  ytick={0,1,2,3,4},yticklabels={{NBIM, gestion interne},{NBIM, fonds entier},{CPP, fonctionnement},{NBIM, mandats externes},{CPP, fonds entier}},
  y dir=reverse,tick label style={font=\footnotesize},grid=major,grid style={gray!20},table/col sep=comma,
  nodes near coords={\pgfmathprintnumber[fixed,precision=1]{\pgfplotspointmeta}},nodes near coords style={font=\scriptsize}]
\addplot[fill=omDef!55,draw=omDef] table[y=k,x=bp]{figdata/industry/08-asset-managers-pensions-sovereign-funds-and-insurers/costs.csv};
\end{axis}
\end{tikzpicture}

\omcaption{Ce que paient deux propriétaires pour gérer leurs actifs, en points de base des actifs de fin d'année, d'après leurs
totaux publiés~: les actifs gérés en interne par le fonds norvégien, l'ensemble du fonds, les charges d'exploitation propres du fonds canadien, les
mandats externes du fonds norvégien (frais de base et commissions de performance), et le total du fonds canadien y compris les frais des
gestionnaires externes. Données~: rapport annuel 2025 et communiqué de résultats de l'exercice 2026, via
\texttt{in\_owners}.}\label{fig:in:asset-managers-pensions-sovereign-funds-and-insurers:costs}
\end{omfigure}

\begin{remark}[Ce que financent les deux modèles]\label{rem:in:asset-managers-pensions-sovereign-funds-and-insurers:models}
Un coût exprimé en points de base n'est que la moitié de la comparaison~: l'autre moitié est ce que ces coûts achètent. Un propriétaire qui verse à des
gestionnaires externes de fortes commissions de performance paie pour des rendements qu'il attend d'eux au-delà de leurs frais, dans des actifs (private
equity, crédit, actifs réels, hedge funds) qu'une équipe indicielle ne peut pas gérer. Un propriétaire dont 95~\% des actifs sont en titres
cotés gérés au plus près d'un indice a besoin de peu de personnes par dollar. La comparaison porte sur des modèles d'affaires, et non
sur l'efficacité.
\end{remark}

Pour un salarié, les deux modèles correspondent à des métiers différents. Le fonds norvégien gère 3.0~milliards de dollars d'actifs par salarié, et
ses frais de personnel s'élèvent à 3.65 millions de NOK par salarié, soit environ 352\,000~dollars au taux moyen de 2025~: moins de personnes, chacune
responsable de montants plus élevés, travaillant surtout sur les marchés cotés. Le modèle du fonds canadien emploie davantage de personnes par dollar
(364 millions de dollars canadiens d'investissements nets par salarié) et beaucoup d'entre eux sélectionnent, supervisent et négocient avec des gestionnaires externes
et des partenaires, ce qui relève d'un tout autre métier que la gestion d'un portefeuille.

\section{Construire ou acheter~: le seuil de rentabilité de l'équipe interne}

\begin{method}[Taille du mandat au seuil de rentabilité]\label{met:in:asset-managers-pensions-sovereign-funds-and-insurers:be}
\begin{enumerate}
\item Chiffrer le coût de l'équipe interne~: effectifs multipliés par le coût du personnel par tête, auxquels s'ajoutent les systèmes et les données.
\item Chiffrer l'option externe~: le taux de frais, en points de base du mandat, auquel s'ajoute la commission de performance attendue.
\item La taille du mandat au seuil de rentabilité est le coût interne divisé par le taux de frais.
\item Au-dessus de ce seuil, l'équipe interne est moins chère, si elle peut obtenir le même rendement~; en dessous, ce sont les frais qui sont plus avantageux.
\end{enumerate}
\end{method}

\begin{example}[Une équipe de dix personnes]\label{ex:in:asset-managers-pensions-sovereign-funds-and-insurers:team}
Une équipe interne de dix personnes à 352\,000~\$ par tête, soit le coût du personnel du fonds norvégien en 2025, auquel s'ajoutent 1.5~millions de dollars de systèmes
et de données (chiffre illustratif), coûte 5.02~millions de dollars par an. Face à des frais externes de 20 points de base, elle atteint le seuil de rentabilité pour un
mandat de 2.51~milliards de dollars~; à 10 points de base, pour 5.02~milliards de dollars~; à 40, pour 1.255~milliards de dollars
(la \cref{fig:in:asset-managers-pensions-sovereign-funds-and-insurers:be}). Les mandats d'un grand propriétaire sont très supérieurs à ces
tailles, ce qui explique que les plus grands propriétaires internalisent les stratégies qu'ils savent gérer.
\end{example}

\begin{omfigure}
\begin{tikzpicture}
\begin{axis}[width=10.5cm,height=5.8cm,xmin=5,xmax=60,ymin=0,ymax=20,xlabel={\small frais externes (points de base du mandat)},
  ylabel={\small mandat au seuil de rentabilité (milliards de dollars)},tick label style={font=\footnotesize},grid=major,grid style={gray!20},
  table/col sep=comma,legend style={legend cell align=left,font=\scriptsize,at={(0.5,-0.3)},anchor=north,/tikz/every even column/.append style={column sep=8pt}},legend columns=3]
\addplot[omDef,thick] table[x=fee_bp,y=team5]{figdata/industry/08-asset-managers-pensions-sovereign-funds-and-insurers/breakeven.csv};
\addplot[omThm,thick,dashed] table[x=fee_bp,y=team10]{figdata/industry/08-asset-managers-pensions-sovereign-funds-and-insurers/breakeven.csv};
\addplot[omProp,thick,dotted] table[x=fee_bp,y=team20]{figdata/industry/08-asset-managers-pensions-sovereign-funds-and-insurers/breakeven.csv};
\legend{5 personnes,10 personnes,20 personnes}
\end{axis}
\end{tikzpicture}

\omcaption{Taille du mandat au-delà de laquelle une équipe interne est moins chère qu'un gérant externe, en fonction des frais externes,
pour des équipes de 5, 10 et 20 personnes à 352\,000~\$ par tête, plus 1.5~millions de dollars de systèmes et de données (illustratif). Données~:
\texttt{firm.ownercost.breakeven\_assets}, via \texttt{in\_owners.team\_breakeven}.}\label{fig:in:asset-managers-pensions-sovereign-funds-and-insurers:be}
\end{omfigure}

\section{Gérants quantitatifs long only}

Les sociétés de gestion qui servent les \omterm{def:in:asset-managers-pensions-sovereign-funds-and-insurers:owner}{propriétaires d'actifs} sont les plus grands employeurs du buy side. Un gérant fait état de 14.0~billions de dollars
d'encours sous gestion et environ 24\,900 salariés, soit quelque 560~millions de dollars par salarié. Le travail quantitatif dans un gérant de
cette taille relève de l'ingénierie d'indices,
la construction de portefeuilles et les produits factoriels (livre~8, chapitre~29), les systèmes de risque et le trading qui
met en œuvre les changements de portefeuille (livre~10, chapitre~20). Le rythme est fixé par les calendriers de rééquilibrage et les mandats des clients, et non par
les mouvements de marché minute par minute, et la rémunération suit les revenus tirés des frais de gestion plutôt que le profit de trading
(chapitre~13).

\section{Assureurs~: gestion actif--passif et quant actuariel}

Un assureur investit les primes en face des engagements qu'il doit à ses assurés, souvent à plusieurs décennies~; sa fonction d'investissement est
celle d'un \omterm{def:in:asset-managers-pensions-sovereign-funds-and-insurers:owner}{propriétaire d'actifs}, et ses quants travaillent sur la gestion actif--passif (livre~6, chapitre~24)~: appariement de la duration et
des flux de trésorerie des actifs aux passifs, couverture du risque de taux d'intérêt et d'inflation, et détention de fonds propres face au décalage
selon les règles de solvabilité. À côté d'eux, les actuaires évaluent eux-mêmes les passifs à l'aide de modèles de mortalité, de résiliation et
de catastrophe. Les deux métiers se rencontrent au bilan, et un quant qui passe de la banque à l'assurance
apporte des modèles de taux et de produits dérivés (livres~2 et~6) et apprend les modèles de passif.

\section{Rythme, rémunération et stabilité côté propriétaires d'actifs}

Le côté \omterm{def:in:asset-managers-pensions-sovereign-funds-and-insurers:owner}{propriétaires d'actifs} offre un compromis différent de celui du reste de ce livre. Les horizons sont longs, le retour d'information est lent et
mesurés par rapport à des indices de référence, et une seule décision --- l'allocation stratégique --- pèse plus que n'importe quelle transaction. La rémunération est fixée par
les conseils d'administration et les gouvernements, publiée de manière agrégée (la ligne des coûts du personnel dans les comptes du fonds norvégien), et rarement liée
à une part du profit. Les institutions durent~: un fonds de retraite ne ferme pas après un drawdown. Pour un quant, ce côté échange
le potentiel de gain contre la stabilité, et l'envergure contre le rythme.

\section{Tutoriel~: construire ou embaucher l'équipe}\label{tut:in:asset-managers-pensions-sovereign-funds-and-insurers:tut}

\begin{tutorial}
\textbf{Objectif.} Mettre sur une même base les coûts publiés de deux \omterm{def:in:asset-managers-pensions-sovereign-funds-and-insurers:owner}{propriétaires d'actifs}, répartir le coût d'un propriétaire entre \omterm{def:in:asset-managers-pensions-sovereign-funds-and-insurers:internal}{gestion interne} et
gestion externe, et calculer la taille du mandat au seuil de rentabilité d'une équipe interne. \textbf{État final~:}
les \cref{fig:in:asset-managers-pensions-sovereign-funds-and-insurers:costs,fig:in:asset-managers-pensions-sovereign-funds-and-insurers:be}.
\begin{tutsteps}
\item \textbf{Les totaux.} \texttt{data/industry/asset\_owners.csv} contient les totaux publiés de chaque propriétaire dans sa propre
monnaie~: actifs, effectifs, coûts du personnel et coûts internes, frais de base et commissions de performance externes, actifs gérés en externe
lorsqu'ils sont publiés.
\item \textbf{Une même base.} \texttt{firm.ownercost.cost\_bp(owner)} divise l'ensemble des coûts de gestion par les actifs de fin d'exercice~;
\texttt{split\_bp} divise les coûts internes par les actifs gérés en interne et les frais externes par les actifs gérés en externe
(le \cref{lst:in:asset-managers-pensions-sovereign-funds-and-insurers:split}).
\omcode{code/firm/ownercost/firm_ownercost.py}{30}{38}{Le coût total et la répartition interne-externe, en points de base.}\label{lst:in:asset-managers-pensions-sovereign-funds-and-insurers:split}
\item \textbf{Par tête.} \texttt{in\_owners.nbim()} convertit les coûts du personnel par salarié en dollars aux taux de change moyens
de la BCE.
\item \textbf{Seuil de rentabilité.} \texttt{in\_owners.team\_breakeven(heads, cost\_per\_head\_k, systems\_k, fees)} applique
la \cref{met:in:asset-managers-pensions-sovereign-funds-and-insurers:be}.
\end{tutsteps}
\end{tutorial}

\textbf{Pour aller plus loin.} Ajouter à l'option externe une commission de performance de 20~\% d'un rendement excédentaire de 1~\%, et recalculer
le seuil de rentabilité (exercice~7)~; ajouter un troisième propriétaire à partir de son rapport annuel, avec ses propres lignes de coûts.

\section{Construction~: le modèle de coût du propriétaire}\label{bld:in:asset-managers-pensions-sovereign-funds-and-insurers:ownercost}

\begin{build}
\bfield{Rôle} Comparer la gestion en interne et l'achat de gestion d'investissement sur une même base, pour ce chapitre et pour les
comparaisons par tête du chapitre~12.
\bfield{Interface} \texttt{firm.ownercost}~:
\texttt{Owner(name, assets, internal\_costs, external\_base, external\_perf, external\_assets)}~; \texttt{cost\_bp}~; \texttt{split\_bp}~; \texttt{internal\_cost}~; \texttt{external\_cost}~;
\texttt{breakeven\_assets}.
\bfield{Règles} Une seule monnaie et un seul dénominateur par comparaison~; la répartition est refusée lorsque le propriétaire ne publie pas ses
actifs gérés en externe~; frais exprimés en points de base du mandat.
\bfield{Tests d'acceptation} \texttt{code/firm/ownercost/tests/}~: coût et répartition sur des totaux construits~; répartition refusée
sans actifs gérés en externe~; coûts internes et externes égaux au seuil de rentabilité.
\bfield{Pour aller plus loin} Coûts de transaction et écart de rendement entre \omterm{def:in:asset-managers-pensions-sovereign-funds-and-insurers:internal}{gestion interne} et gestion externe~; une équipe dont
le coût par tête augmente avec la complexité du mandat.
\end{build}

\begin{omsources}
\item Norges Bank Investment Management, Government Pension Fund Global, rapport annuel 2025 (tableau 12.1).
\item CPP Investments, communiqué de résultats de l'exercice 2026 (20 mai 2026).
\item International Working Group of Sovereign Wealth Funds, Santiago Principles (2008)~; US Internal Revenue Service,
defined benefit plans.
\item BlackRock, Form 10-K for 2025~; Banque centrale européenne, taux de référence annuels moyens.
\end{omsources}

\section{Exercices}

\begin{exercise}[$\star$]\label{exo:in:asset-managers-pensions-sovereign-funds-and-insurers:1}
Convertir la valeur du fonds norvégien en dollars américains aux taux moyens de 2025 de la BCE, et calculer ses actifs par salarié.
\end{exercise}

\begin{exercise}[$\star$]\label{exo:in:asset-managers-pensions-sovereign-funds-and-insurers:2}
Calculer les frais de base versés par le fonds norvégien aux gérants externes, en points de base des actifs gérés en externe.
\end{exercise}

\begin{exercise}[$\star$]\label{exo:in:asset-managers-pensions-sovereign-funds-and-insurers:3}
Un régime de retraite promet une prestation fixe au départ à la retraite. Qui porte le risque d'investissement, et quel travail quantitatif en découle~?
\end{exercise}

\begin{exercise}[$\star\star$]\label{exo:in:asset-managers-pensions-sovereign-funds-and-insurers:4}
Calculer les frais versés par le fonds canadien à ses gérants externes, ainsi que ses frais de fonctionnement propres, en points de base des actifs nets de fin d'année,
et expliquer pourquoi son ratio de frais de fonctionnement publié (23.1) diffère du chiffre de fonctionnement obtenu.
\end{exercise}

\begin{exercise}[$\star\star$]\label{exo:in:asset-managers-pensions-sovereign-funds-and-insurers:5}
Une équipe interne de 20 personnes coûte 352\,000~\$ par tête, plus 1.5~millions de dollars de systèmes. Quelle taille de mandat atteint le point d'équivalence
face à des frais externes de 25 points de base~?
\end{exercise}

\begin{exercise}[$\star\star$]\label{exo:in:asset-managers-pensions-sovereign-funds-and-insurers:6}
Lire la \cref{fig:in:asset-managers-pensions-sovereign-funds-and-insurers:be}~: pour une équipe de dix personnes, quels frais rendent un mandat
de 2~milliards de dollars équivalent~?
\end{exercise}

\begin{exercise}[$\star\star\star$]\label{exo:in:asset-managers-pensions-sovereign-funds-and-insurers:7}
\emph{Programmation.} Un gérant externe facture 20 points de base, plus 20~\% d'une surperformance attendue de 1~\% par an. Recalculer
le mandat d'équivalence de l'équipe de dix personnes. Que doit livrer l'équipe interne pour que la comparaison tienne~?
\end{exercise}

\begin{exercise}[$\star\star\star$]\label{exo:in:asset-managers-pensions-sovereign-funds-and-insurers:8}
\emph{Trouver la faille.} «~Le fonds norvégien coûte 3.5 points de base et le fonds canadien 81.8~: le fonds canadien est 23 fois
moins efficace.~»
\end{exercise}

\section{Problème~: construire ou embaucher l'équipe}

\begin{problem}[{Problème du week-end --- construire ou embaucher l'équipe}]\label{pb:in:asset-managers-pensions-sovereign-funds-and-insurers:1}
Le directeur des investissements d'un fonds de pension disposant de 200~milliards de dollars d'actifs demande
s'il faut rapatrier en interne un mandat actions de 15~milliards de dollars, et un quant doit répondre à partir de chiffres publics.

\medskip\noindent\textbf{Partie I --- Les propriétaires.}
\begin{enumerate}
\item Définir un \omterm{def:in:asset-managers-pensions-sovereign-funds-and-insurers:owner}{propriétaire d'actifs}, un \omterm{def:in:asset-managers-pensions-sovereign-funds-and-insurers:swf}{fonds souverain} et un \omterm{def:in:asset-managers-pensions-sovereign-funds-and-insurers:swf}{régime de retraite à prestations définies}.
\item Définir la \omterm{def:in:asset-managers-pensions-sovereign-funds-and-insurers:internal}{gestion interne} et donner la part interne du fonds norvégien.
\item Citer les cinq types de travail quantitatif chez un \omterm{def:in:asset-managers-pensions-sovereign-funds-and-insurers:owner}{propriétaire d'actifs}.
\item Donner les ratios de coûts publiés des deux propriétaires, et dire pourquoi ils ne peuvent pas être comparés.
\item En quoi consistent les frais de gestion du fonds norvégien~?
\end{enumerate}

\medskip\noindent\textbf{Partie II --- Un même périmètre.}
\begin{enumerate}[resume]
\item Calculer les coûts de gestion totaux des deux propriétaires, en points de base des actifs de fin d'année.
\item Répartir le coût du fonds norvégien entre \omterm{def:in:asset-managers-pensions-sovereign-funds-and-insurers:internal}{gestion interne} et gestion externe.
\item Quelle part du coût total du fonds canadien représentent les frais des gérants externes~?
\item Donner le coût du personnel par salarié du fonds norvégien, en couronnes puis en dollars.
\item Donner ses actifs par salarié, en dollars.
\end{enumerate}

\medskip\noindent\textbf{Partie III --- L'équipe.}
\begin{enumerate}[resume]
\item Énoncer la méthode du point d'équivalence.
\item Chiffrer une équipe de dix personnes au coût du personnel par tête du fonds norvégien et 1.5~millions de dollars de systèmes.
\item Donner son mandat d'équivalence à des frais externes de 10, 20 et 40 points de base.
\item À 20 points de base, combien de fois le mandat d'équivalence représente-t-il le mandat de 15~milliards de dollars~?
\item Combien coûterait le gérant externe sur 15~milliards de dollars à 20 points de base~?
\end{enumerate}

\medskip\noindent\textbf{Partie IV --- Le verdict.}
\begin{enumerate}[resume]
\item Énoncer le \emph{résultat nommé}~: les tailles de mandat d'équivalence d'une équipe de dix personnes à 10, 20 et 40 points de base.
\item Que doit livrer l'équipe interne pour que la réponse tienne~?
\item Qu'est-ce que la comparaison des coûts des deux propriétaires laisse de côté~?
\item Quels emplois la reprise du mandat en interne crée-t-elle, et lesquels supprime-t-elle~?
\item En deux phrases, que doit dire le quant au directeur des investissements~?
\end{enumerate}
\end{problem}

\section{Questions d'entretien}

\begin{interviewq}[$\star$ \iqroles{chercheur}]\label{iq:in:asset-managers-pensions-sovereign-funds-and-insurers:1}
Quelle est la différence entre un \omterm{def:in:asset-managers-pensions-sovereign-funds-and-insurers:owner}{propriétaire d'actifs} et une société de gestion, et comment change-t-elle le métier d'un quant~?
\end{interviewq}

\begin{interviewq}[$\star$ \iqroles{risque}]\label{iq:in:asset-managers-pensions-sovereign-funds-and-insurers:2}
La duration du passif d'un régime de retraite est de 18 ans, et celle de ses obligations de 6 ans. Que devient son
ratio de financement quand les taux longs baissent, et que feriez-vous pour y remédier~?
\end{interviewq}

\begin{interviewq}[$\star\star$ \iqroles{chercheur}]\label{iq:in:asset-managers-pensions-sovereign-funds-and-insurers:3}
Un propriétaire verse à un gérant externe 50 points de base, plus 20~\% de la surperformance. Quel rendement excédentaire le gérant doit-il
espérer dégager pour justifier son embauche face à un portefeuille indiciel interne de 5 points de base~?
\end{interviewq}

\begin{interviewq}[$\star\star$ \iqroles{chercheur, risque}]\label{iq:in:asset-managers-pensions-sovereign-funds-and-insurers:4}
Pourquoi l'allocation stratégique d'actifs domine-t-elle le rendement d'un grand fonds, et qu'est-ce que cela implique
pour l'endroit où ses meilleurs quants doivent travailler~?
\end{interviewq}

\begin{interviewq}[$\star\star$ \iqroles{banque, risque}]\label{iq:in:asset-managers-pensions-sovereign-funds-and-insurers:5}
Lesquelles de vos compétences de quant taux en banque transfèrent-elles à l'équipe actif--passif d'un assureur, et que devriez-vous
apprendre~?
\end{interviewq}

\begin{interviewq}[$\star\star\star$ \iqroles{chercheur}]\label{iq:in:asset-managers-pensions-sovereign-funds-and-insurers:6}
Concevez un test permettant de savoir si l'équipe actions interne d'un propriétaire crée de la valeur nette de ses coûts, à partir de dix ans de rendements mensuels
de l'équipe et de son indice de référence.
\end{interviewq}
